% ============================================================
%  Prolog
% ============================================================
\chapter*{Prolog}
\addcontentsline{toc}{chapter}{Prolog}
\markboth{Prolog}{Prolog}

Dieses Skript ist meine Begleitdokument zur Vorlesung
\textbf{Systemtheorie – Signale und Systeme~I} im Sommersemester~2026.
Es soll als mathematisch strukturierte Mitschrift mit dem Ziel, die
zentralen Konzepte und Methoden der Signalverarbeitung und Systemanalyse
kompakt und nachvollziehbar festzuhalten, geschrieben werden. 
Der Inhalt soll mathematisch komplexe sProzesse verständlich erklären.

\bigskip

\begin{bemerkungbox}[Hinweis zur Nutzung]
Das Skript umfasst mehr Inhalte als die Vorlesung da Grundlagen mit erklärt 
werden.
\end{bemerkungbox}

\bigskip

\section*{Aufbau des Skripts}

Das Skript folgt dem Verlauf der Vorlesung und ist in drei Teile gegliedert:

\begin{enumerate}
  \item \textbf{Vorlesungen} – Kapitelweise Mitschrift der einzelnen
        Vorlesungseinheiten in chronologischer Reihenfolge.
  \item \textbf{Zusätzliches Wissen} – Themen, die über den direkten
        Vorlesungsinhalt hinausgehen, aber für das Verständnis relevant
        oder mathematisch interessant sind.
  \item \textbf{Formelsammlung} – Eine kompakte Zusammenstellung aller
        wichtigen Formeln und Identitäten zum schnellen Nachschlagen.
\end{enumerate}

\section*{Verwendete Boxtypen}

Zur übersichtlichen Darstellung unterschiedlicher Inhaltstypen werden
folgende farblich kodierte Boxen verwendet:

\begin{definitionbox}[Beispiel]
  Formale Definitionen mathematischer Begriffe und Konzepte.
\end{definitionbox}

\begin{satzbox}[Beispiel]
  Mathematische Sätze und Theoreme mit Anspruch auf allgemeine Gültigkeit.
\end{satzbox}

\begin{beweisbox}[Beispiel]
  Beweise zu Sätzen, sofern sie für das Verständnis relevant sind.
\end{beweisbox}

\begin{exercisebox}[Beispiel]
  Anschauliche Beispiele und Rechenaufgaben zur Illustration der Theorie.
\end{exercisebox}

\begin{bemerkungbox}[Beispiel]
  Ergänzende Hinweise, Besonderheiten und Querverweise.
\end{bemerkungbox}

\begin{gedankenbox}[Beispiel]
  Gedankenexperimente zur intuitiven Motivation von Konzepten.
\end{gedankenbox}

\vfill
\begin{flushright}
  \textit{Philipp Lehmann, Sommersemester 2026}
\end{flushright}
